\chapter{Fat-soluble vitamins: A, D, E, and K}

\section{Learning objectives}
Describe the physiological roles and food forms of vitamins A, D, E, and K; recognize why fat malabsorption increases risk; distinguish deficiency syndromes from nonspecific symptoms; and explain why chronic high-dose use is clinically different from food intake.

\section{Vitamin A and carotenoids}
Vitamin A supports vision, immune defenses, epithelial tissues, growth, and reproduction. Preformed vitamin A (retinol and retinyl esters) occurs in liver, fish, eggs, and dairy. Orange, yellow, and dark-green plants supply provitamin A carotenoids such as beta-carotene; conversion varies between people. Severe deficiency can impair night vision and immunity, while excess preformed vitamin A can harm the liver, bone, and a developing fetus. Beta-carotene supplements are not a safe substitute for food and are especially concerning for smokers and former smokers. \cite{odsA}

Retinal participates in the visual cycle, while retinoic acid regulates gene transcription and epithelial differentiation. Deficiency first affects dark adaptation and mucosal integrity; xerophthalmia and keratomalacia are severe ocular manifestations. Diagnosis may involve dietary history, examination, and serum retinol, although serum concentrations are homeostatically controlled and can fall during infection. In pregnancy, preformed retinoids require particular caution because teratogenicity is dose- and timing-dependent. Carotenoid-rich foods can discolor skin harmlessly, but this is not the same as vitamin A toxicity.

\section{Vitamin D}
Vitamin D helps regulate calcium and phosphate and is important for bone mineralization. Skin can produce vitamin D with ultraviolet B exposure, but season, latitude, skin pigmentation, age, clothing, sunscreen, and indoor living affect production. Food sources include fatty fish, egg yolk, fortified milk or plant beverages, and some mushrooms. Deficiency can cause rickets in children and osteomalacia in adults. High-dose self-treatment can produce hypercalcemia and kidney injury; more vitamin D has not been shown to prevent every chronic disease. \cite{odsD}

Vitamin D status is commonly assessed with serum 25-hydroxyvitamin D; the active hormone 1,25-dihydroxyvitamin D is regulated differently and is not a routine status test. Interpretation depends on the assay, laboratory, clinical condition, and guideline used. Risk rises with limited sun exposure, dark skin at high latitude, malabsorption, obesity, liver or kidney disease, and some anticonvulsants. Bone health also requires mechanical loading, adequate calcium and protein, and evaluation for secondary causes of osteoporosis.

\section{Vitamin E}
Vitamin E is a family of compounds with antioxidant activity and roles in cell signaling and immune function. Nuts, seeds, vegetable oils, wheat germ, and some green vegetables are useful sources. Clinically important deficiency is uncommon but can occur with fat-malabsorption disorders. High-dose supplements can increase bleeding risk and may interact with anticoagulants; routine high-dose use for prevention is not supported.

Alpha-tocopherol is the form preferentially maintained in human plasma. Deficiency may cause peripheral neuropathy, ataxia, skeletal-myopathy, or retinal injury, usually in the setting of severe malabsorption or genetic disorders rather than ordinary dietary variation. Treatment is directed at the underlying disorder and may require specialist dosing. Antioxidant theory alone is insufficient evidence for supplementation: redox biology is complex, and high doses can disrupt signaling or interact with treatment.

\section{Vitamin K}
Vitamin K is needed for blood-clotting proteins and bone-related proteins. Leafy greens are rich in phylloquinone (K1); intestinal bacteria and animal foods contribute other forms. Consistency matters for people taking warfarin: do not abruptly eliminate greens or begin a K supplement without medical advice. Newborns receive vitamin K prophylaxis in many health systems because stores and transfer across the placenta are limited.

Vitamin K deficiency is unusual in healthy adults but can occur with prolonged antibiotic exposure, cholestasis, severe fat malabsorption, or inadequate neonatal stores. Prolonged clotting time may be an early clue, although many causes of coagulopathy must be considered. Vitamin K is not a general clotting enhancer and should not be self-administered to reverse anticoagulation.

\section{Malabsorption and clinical integration}
Because these vitamins depend on bile and intestinal lipid handling, chronic pancreatitis, cholestatic disease, inflammatory bowel disease, bariatric surgery, and some genetic disorders can produce combined deficiencies. Chronic diarrhea, steatorrhea, unexplained weight loss, bruising, visual changes, or bone pain merits medical evaluation. Management may require dietary modification, treatment of the underlying disease, and monitored replacement rather than an over-the-counter multivitamin.

\section{Practical summary}
Use colorful plants, fortified foods, and modest amounts of healthy fats. Avoid cod-liver-oil or multi-supplement stacking without checking total vitamin A and D.
